% chapters/01-introduction.tex — บทที่ 1 บทนำ (expanded)
\chapter{บทนำ}
\label{ch:intro}

\section{ที่มาและความสำคัญของงานวิจัย}
\label{sec:background}

อุตสาหกรรมก่อสร้างเป็นหนึ่งในภาคเศรษฐกิจที่มีบทบาทสำคัญที่สุดต่อการพัฒนา
ประเทศ ทั้งในแง่การสร้างงาน การลงทุนภาครัฐ และการพัฒนาโครงสร้างพื้นฐาน
ที่จำเป็นต่อคุณภาพชีวิตของประชาชน \cite{sirieawphikul2024}. โครงการก่อสร้าง
ภาครัฐในประเทศไทยมีมูลค่ารวมหลายแสนล้านบาทต่อปี ครอบคลุมตั้งแต่งานถนน
สะพาน อาคารสำนักงาน โรงพยาบาล โรงเรียน และระบบสาธารณูปโภคพื้นฐาน
ความถูกต้องในการประมาณราคาโครงการเหล่านี้จึงเป็นปัจจัยสำคัญที่ส่งผลโดยตรง
ต่อความคุ้มค่าของงบประมาณและความโปร่งใสในการจัดซื้อจัดจ้างของหน่วยงานภาครัฐ.

หัวใจสำคัญของการประมาณราคางานก่อสร้างคือการกำหนดราคาต่อหน่วย (Unit Cost)
ของวัสดุก่อสร้าง ซึ่งเป็นองค์ประกอบหลักของต้นทุนงาน \cite{ashworth2013}.
ในประเทศที่มีระบบประมาณราคามาตรฐาน ผู้ประมาณราคาจะใช้ฐานข้อมูลราคาวัสดุ
ที่หน่วยงานกลางเผยแพร่เป็นจุดอ้างอิง อย่างไรก็ตาม ความถูกต้องของการประมาณ
ขึ้นอยู่กับสามปัจจัยหลัก ได้แก่ (1) ความครบถ้วนของข้อมูลราคา (2) ความทันสมัย
ของข้อมูล และ (3) ความสอดคล้องของราคามาตรฐานกับราคาตลาดจริง
\cite{crone1992,leung2005}.

ในบริบทประเทศไทย หน่วยงานหลักที่เผยแพร่ราคามาตรฐานวัสดุก่อสร้าง ได้แก่
สำนักงานนโยบายและยุทธศาสตร์การค้า (สนค.) ของกระทรวงพาณิชย์ ซึ่งเผยแพร่
ดัชนีราคาวัสดุก่อสร้าง (Construction Materials Index: CMI) เป็นรายเดือน
\cite{tpso2569}, และกรมบัญชีกลาง (CGD) ที่กำหนดราคามาตรฐานสำหรับงาน
ก่อสร้างของส่วนราชการตามหลักเกณฑ์การคำนวณราคากลางงานก่อสร้าง
\cite{cgd2569}. ทั้งสองแหล่งให้ข้อมูลที่เชื่อถือได้และเป็นทางการ แต่มี
ข้อจำกัดสำคัญหลายประการ.

ประการแรก ข้อมูลราคาทางการมีรอบการปรับปรุงค่อนข้างยาว (รายเดือนหรือ
รายไตรมาส) ทำให้ในช่วงที่ราคาตลาดผันผวน ข้อมูลที่ใช้อ้างอิงอาจไม่สะท้อน
ราคาที่แท้จริง ประการที่สอง ราคาทางการมักเป็นราคาเฉลี่ยระดับประเทศหรือ
ระดับภาค ไม่ครอบคลุมความหลากหลายเชิงพื้นที่ของจังหวัดต่าง ๆ ที่อาจมีต้นทุน
ขนส่ง อุปสงค์-อุปทาน และโครงสร้างตลาดท้องถิ่นแตกต่างกัน \cite{leung2005}.
ประการที่สาม รายการวัสดุที่ภาครัฐเผยแพร่อาจไม่ครอบคลุมวัสดุเฉพาะทางหรือ
ยี่ห้อเฉพาะที่ใช้ในโครงการจริง ทำให้ผู้ประมาณราคาต้องสืบค้นข้อมูลเพิ่มเติม
จากแหล่งอื่น \cite{benchanakatkun2014}.

ในขณะเดียวกัน ภาคค้าปลีกวัสดุก่อสร้างสมัยใหม่ (Modern Trade) ได้พัฒนาช่องทาง
การจำหน่ายออนไลน์อย่างกว้างขวาง ผู้ประกอบการรายใหญ่ เช่น HomePro,
MegaHome, บุญถาวร, ไทวัสดุ, โกลบอลเฮ้าส์, ดูโฮม, เอสซีจีโฮม และ บีเอ็นบีโฮม
ได้เผยแพร่ราคาสินค้าผ่านเว็บไซต์ของตน พร้อมข้อมูลรายละเอียดสินค้า ขนาด
คุณสมบัติ และยี่ห้อ การปรับปรุงราคาเหล่านี้ทำในรอบที่สั้นกว่าราคาทางการ
จึงสามารถสะท้อนราคาตลาดจริงได้ดียิ่งกว่า \cite{elmohr2022}.

หากสามารถนำข้อมูลราคาจากแหล่งภาครัฐและภาคเอกชนมาประมวลผลร่วมกันอย่าง
เป็นระบบ จะช่วยเพิ่มทั้งความครบถ้วนและความแม่นยำของข้อมูลราคาวัสดุ
ก่อสร้างที่ใช้ในการประมาณราคา อย่างไรก็ตาม กระบวนการรวบรวมข้อมูลจาก
หลายแหล่งในปัจจุบันยังเป็นแบบกึ่งอัตโนมัติหรือทำด้วยมือ ใช้เวลานานและ
ขาดความเป็นมาตรฐานเดียวกัน \cite{benchanakatkun2014}. นอกจากนี้ ข้อมูล
ราคาจากแหล่งต่าง ๆ มีรูปแบบที่ไม่เป็นมาตรฐาน — บ้างเผยแพร่เป็นเอกสาร PDF
บ้างเป็นตาราง HTML บ้างซ่อนอยู่ในเว็บแอปพลิเคชันแบบ Single Page Application
(SPA) — ทำให้การจัดทำฐานข้อมูลรวมเป็นภาระงานที่ซับซ้อน
\cite{rao2015,chunyaem2019}.

ปัญหาดังกล่าวจึงเป็นช่องว่างของงานวิจัยที่งานวิจัยนี้มุ่งจะตอบสนอง โดยพัฒนา
ระบบอัตโนมัติเพื่อรวบรวมข้อมูลราคาวัสดุก่อสร้างจากแหล่งหลายแหล่ง
ทั้งภาครัฐและภาคเอกชน ผ่านเทคนิค Web Scraping, Application Programming
Interface (API), และระบบบูรณาการฐานข้อมูล (Database Integration) ลดการพึ่งพา
แรงงานคน ลดความคลาดเคลื่อนของข้อมูล และเพิ่มความสามารถในการเปรียบเทียบ
ราคาในเชิงพื้นที่และเชิงเวลา ซึ่งเป็นข้อมูลพื้นฐานที่จำเป็นต่อการคำนวณ
ต้นทุนต่อหน่วยและการจัดทำบัญชีแสดงปริมาณงาน (Bill of Quantities: BOQ)
ของโครงการก่อสร้าง.

ระบบที่พัฒนาขึ้นสามารถนำไปประยุกต์ใช้ในหลายบริบท ได้แก่ การประมาณราคา
โครงการของหน่วยงานภาครัฐที่ต้องการอ้างอิงราคามาตรฐานควบคู่กับราคาตลาด
จริง การวิเคราะห์ความผันผวนของราคาวัสดุสำหรับใช้วางแผนงบประมาณ
และการสร้างความโปร่งใสในกระบวนการจัดซื้อจัดจ้างภาครัฐผ่านการเปิดเผย
ข้อมูลราคาจากหลายแหล่งให้ผู้มีส่วนได้ส่วนเสียสามารถตรวจสอบเปรียบเทียบ
ได้อย่างเป็นระบบ.

\section{ปัญหางานวิจัย}
\label{sec:problem}

จากการศึกษาเบื้องต้นพบว่าปัญหาในกระบวนการสืบค้นและคำนวณต้นทุนวัสดุ
ก่อสร้างในปัจจุบันสามารถสรุปได้ดังนี้

\begin{enumerate}[leftmargin=2em,itemsep=0.4em]
  \item \textbf{ความกระจัดกระจายของแหล่งข้อมูลราคา}\\
        ราคาวัสดุก่อสร้างเผยแพร่ในหลายแหล่งที่ไม่เชื่อมโยงกัน ทั้งเว็บไซต์
        ผู้จำหน่ายปลีก ร้านค้าวัสดุ Modern Trade เอกสารหน่วยงานภาครัฐ
        และฐานข้อมูลเฉพาะองค์กร ผู้ประมาณราคาต้องสืบค้นจากแต่ละแหล่ง
        แล้วนำมารวมด้วยตนเอง ซึ่งใช้เวลานานและมีโอกาสเกิดความคลาดเคลื่อน
        จากการป้อนข้อมูลด้วยมือ \cite{rao2015}.

  \item \textbf{ความไม่เป็นมาตรฐานของรูปแบบข้อมูล}\\
        ข้อมูลราคาจากแหล่งต่าง ๆ มีโครงสร้างหลากหลาย ทั้งหน่วยจำหน่าย
        (เช่น บาทต่อถุง 50 กก., บาทต่อตัน, บาทต่อเส้น) ชื่อเรียกสินค้า
        ขนาด เกรด มาตรฐานการผลิต และรูปแบบเอกสาร (PDF, HTML, JSON)
        จึงต้องมีกระบวนการแปลงหน่วย ทำความสะอาดข้อมูล และจับคู่รายการ
        ก่อนนำมาประมวลผล \cite{chunyaem2019}.

  \item \textbf{ความไม่ทันสมัยของข้อมูลราคา}\\
        ราคาวัสดุเปลี่ยนแปลงตามภาวะตลาด ต้นทุนพลังงาน และค่าขนส่ง
        ในขณะที่ราคามาตรฐานภาครัฐมีรอบการเผยแพร่รายเดือนหรือรายไตรมาส
        \cite{tpso2569}. ระยะห่างนี้อาจส่งผลให้ราคาที่นำมาใช้ไม่สะท้อน
        ราคาตลาดปัจจุบัน โดยเฉพาะในช่วงราคาวัสดุก่อสร้างผันผวนสูง.

  \item \textbf{ความแตกต่างของราคาตามพื้นที่}\\
        ราคาวัสดุก่อสร้างมีความแตกต่างตามภูมิภาค ระยะทางขนส่งจากแหล่งผลิต
        และโครงสร้างตลาดท้องถิ่น แต่แหล่งข้อมูลราคาส่วนใหญ่ไม่ระบุ
        พื้นที่จัดจำหน่ายอย่างชัดเจน หรือให้เพียงราคาเฉลี่ยระดับประเทศ
        ทำให้การนำไปใช้ในพื้นที่ก่อสร้างเฉพาะอาจไม่สอดคล้องกับสภาพจริง.

  \item \textbf{ข้อจำกัดในการเข้าถึงข้อมูลดิจิทัล}\\
        ข้อมูลบางส่วนอยู่ในรูปแบบ PDF หรือรูปภาพที่ไม่สามารถดึงข้อมูล
        เชิงโครงสร้างได้โดยตรง บางเว็บไซต์ใช้เทคโนโลยี Single Page
        Application (SPA) ที่โหลดข้อมูลผ่านการเรียก JavaScript หลัง
        hydrate ทำให้การ scrape ด้วยวิธีดั้งเดิมไม่สามารถดึงราคาได้
        \cite{rao2015}. นอกจากนี้บางเว็บไซต์มีระบบป้องกันบอท เช่น
        Cloudflare bot challenge หรือ CAPTCHA ที่บล็อกการเข้าถึงอัตโนมัติ.

  \item \textbf{ความไม่สอดคล้องของชื่อและคุณลักษณะวัสดุข้ามแหล่ง}\\
        วัสดุก่อสร้างชนิดเดียวกันอาจมีชื่อเรียก ขนาด หรือเกรดที่แตกต่างกัน
        ในแต่ละแหล่งข้อมูล เช่น ปูนซีเมนต์ Type I อาจถูกขายเป็น
        ``ปูนตราเสือ 50 กก.'' ใน HomePro แต่เป็น ``Tiger Cement Type I 50kg''
        ในเว็บอื่น ทำให้การจับคู่รายการ (entity matching) ทำได้ยาก
        และอาจเลือกใช้ราคาที่ไม่ตรงกับวัสดุจริง \cite{chunyaem2019}.

  \item \textbf{การคำนวณต้นทุนต่อหน่วยที่ยังไม่บูรณาการ}\\
        แม้จะรวบรวมราคาได้แล้ว การคำนวณต้นทุนต่อหน่วยตามหลักเกณฑ์
        การคำนวณราคากลางของภาครัฐยังต้องทำด้วยมือ \cite{cgd2569} ผู้ประมาณ
        ราคาต้องคำนวณปริมาณการใช้วัสดุต่อหน่วยงาน คูณกับราคาต่อหน่วย
        แล้วรวมเป็นต้นทุนรวม กระบวนการนี้ใช้เวลามากเมื่อปริมาณงานสูง
        และเปิดโอกาสให้เกิดความคลาดเคลื่อน.

  \item \textbf{การขาดเครื่องมือเปรียบเทียบและวิเคราะห์เชิงสถิติ}\\
        การวิเคราะห์ความสอดคล้องระหว่างราคาทางการกับราคาตลาด การหา
        ค่าผิดปกติ (outlier) และการวิเคราะห์แนวโน้มราคา ยังขาดเครื่องมือ
        ที่ใช้งานง่ายและบูรณาการกับฐานข้อมูลราคา \cite{lee2024,liu2024}.
\end{enumerate}

\section{วัตถุประสงค์ของงานวิจัย}
\label{sec:objectives}

งานวิจัยนี้มีวัตถุประสงค์หลัก 3 ประการ ดังนี้

\begin{enumerate}[leftmargin=2em,itemsep=0.4em]
  \item \textbf{เพื่อพัฒนาระบบอัตโนมัติในการดึงข้อมูลราคาวัสดุก่อสร้าง}
        จากเว็บไซต์ของกระทรวงพาณิชย์ (TPSO, CGD, DIT) และจากเว็บไซต์
        ของผู้ค้าปลีกวัสดุก่อสร้างสมัยใหม่ (Modern Trade) จำนวน 8 แห่ง
        ครอบคลุมทุกจังหวัดของประเทศไทย โดยอาศัยเทคนิค Web Scraping,
        JSON API consumption, และระบบ ScrapingBee proxy เพื่อรองรับเว็บไซต์
        ที่มีระบบป้องกันบอท พร้อมระบบ fallback ผ่าน Cloudflare Browser
        Rendering สำหรับเว็บไซต์ที่ render ด้วย JavaScript.

  \item \textbf{เพื่อพัฒนาระบบการคำนวณต้นทุนวัสดุต่อหน่วย}
        ตามหลักเกณฑ์การคำนวณราคากลางของภาครัฐ โดยรองรับ 5 ประเภทงานหลัก
        ได้แก่ งานบุกระเบื้องผนัง งานเสาเอ็น/คานทับหลัง งานเหล็กเสริม
        งานทาสี และงานก่ออิฐ ระบบสามารถดึงราคาวัสดุจากแหล่งที่ผู้ใช้เลือก
        และคำนวณต้นทุนรวมตามปริมาณงานที่กรอก พร้อมแสดงรายการ Bill of
        Materials (BOM) แบบละเอียดที่สามารถ export เป็น Excel/PDF เพื่อใช้
        ในเอกสารโครงการ.

  \item \textbf{เพื่อวิเคราะห์แนวโน้มและความสอดคล้องของต้นทุนวัสดุต่อหน่วย}
        ในแต่ละประเภทตามช่วงเวลาและพื้นที่ โดยพัฒนาเครื่องมือวิเคราะห์
        เชิงสถิติพื้นฐาน ได้แก่ การวิเคราะห์ค่าเฉลี่ย (Mean Analysis) การคำนวณ
        อัตราการเปลี่ยนแปลงร้อยละ (Percentage Change) การวิเคราะห์แนวโน้ม
        เชิงเส้น (Linear Trend) และการตรวจจับค่าผิดปกติ (Outlier Detection)
        ด้วย z-score test \cite{lee2024} เพื่อให้ระบบสามารถสนับสนุนการตัดสินใจ
        เชิงราคาที่ใช้ในงานก่อสร้างได้อย่างน่าเชื่อถือ.
\end{enumerate}

\section{ขอบเขตของการศึกษา}
\label{sec:scope}

\subsection{ขอบเขตด้านแหล่งข้อมูล}

\begin{itemize}[leftmargin=2em]
  \item \textbf{แหล่งข้อมูลภาครัฐ (3 แหล่ง):}
        \begin{itemize}
          \item TPSO (สนค.) — ดัชนี CMI รายจังหวัด/หมวดวัสดุ ย้อนหลัง 9 เดือน
          \item CGD (กรมบัญชีกลาง) — ราคามาตรฐานสำหรับงานก่อสร้างราชการ
          \item DIT (กรมการค้าภายใน) — ราคาตลาดรายวัน (เมื่อเข้าถึงได้)
        \end{itemize}
  \item \textbf{แหล่งข้อมูลภาคเอกชน (8 แหล่ง):}
        HomePro, MegaHome, บุญถาวร (Boonthavorn), ไทวัสดุ (Thai Watsadu),
        โกลบอลเฮ้าส์ (Global House), ดูโฮม (Dohome), เอสซีจีโฮม (SCG Home),
        และ บีเอ็นบีโฮม (BnB Home).
\end{itemize}

\subsection{ขอบเขตด้านวัสดุก่อสร้าง}

ครอบคลุมวัสดุก่อสร้างหลัก 24 รายการ จัดเป็น 6 หมวด ได้แก่
(1) กระเบื้องและวัสดุปูพื้น/ผนัง (5 รายการ),
(2) ปูนซีเมนต์/ทราย/หิน (4 รายการ),
(3) เหล็กเสริม (8 รายการ ตั้งแต่ RB6 ถึง DB25),
(4) สี (3 รายการ — ทาภายใน/ภายนอก/รองพื้น),
(5) อิฐ (2 รายการ — อิฐมอญและอิฐมวลเบา), และ
(6) คอนกรีตผสมเสร็จ (2 รายการ — กำลัง 240 ksc และ 280 ksc).

\subsection{ขอบเขตด้านพื้นที่}

ระบบรองรับการแยกข้อมูลรายจังหวัดทั้ง 77 จังหวัดของประเทศไทย โดยมีจังหวัด
อ้างอิงเริ่มต้นเป็นกรุงเทพมหานคร (province ID = 10) และมีการเก็บราคา
เปรียบเทียบ 6 ภูมิภาค ได้แก่ ภาคกลาง ภาคเหนือ ภาคตะวันออกเฉียงเหนือ
ภาคตะวันออก ภาคใต้ และภาคตะวันตก.

\subsection{ขอบเขตด้านช่วงเวลา}

ข้อมูลย้อนหลังสำหรับการวิเคราะห์แนวโน้ม ครอบคลุมช่วงเวลา 9 เดือน
ตั้งแต่เดือนสิงหาคม พ.ศ.\ 2568 ถึงเดือนเมษายน พ.ศ.\ 2569 อ้างอิงจาก
ดัชนี TPSO CMI \cite{tpso2569}.

\subsection{ขอบเขตด้านวิธีการ}

ใช้เทคนิค Web Scraping ผ่านเครื่องมือ \texttt{@cloudflare/puppeteer},
ScrapingBee API, และ \texttt{unpdf} สำหรับการสกัดข้อมูลจาก PDF
ระบบ deploy บน Cloudflare Workers (edge runtime) โดยใช้ Cloudflare KV
เป็น storage หลัก ขั้นตอนการวิเคราะห์ใช้สถิติพื้นฐาน — ค่าเฉลี่ย
ค่ามัธยฐาน ส่วนเบี่ยงเบนมาตรฐาน z-score และการวิเคราะห์เส้นแนวโน้ม
แบบ least-squares regression.

\section{ประโยชน์ที่คาดว่าจะได้รับ}
\label{sec:benefit}

\begin{enumerate}[leftmargin=2em]
  \item \textbf{ด้านวิชาการ:} ได้ระบบต้นแบบที่บูรณาการเทคนิค Web Scraping
        + canonicalization + cost calculation + statistical analysis
        เข้าด้วยกัน ซึ่งเติมเต็มช่องว่างของงานวิจัยก่อนหน้าที่มักศึกษา
        เพียงส่วนใดส่วนหนึ่ง \cite{rao2015,chunyaem2019,liu2024}.
  \item \textbf{ด้านการประยุกต์ใช้:} ผู้ประมาณราคา สถาปนิก วิศวกรโยธา
        และเจ้าหน้าที่ภาครัฐสามารถใช้ระบบในการประมาณราคาที่อ้างอิงได้
        จากหลายแหล่งพร้อมกัน ลดเวลาและความคลาดเคลื่อน.
  \item \textbf{ด้านความโปร่งใสและธรรมาภิบาล:} การเปรียบเทียบราคา
        ทางการกับราคาตลาดอย่างเป็นระบบช่วยเพิ่มความโปร่งใสในการจัดซื้อ
        จัดจ้างภาครัฐ และเป็นข้อมูลสนับสนุนการตัดสินใจที่ตรวจสอบได้.
  \item \textbf{ด้านข้อมูลเปิด (Open Data):} ระบบให้บริการ public REST API
        ที่ผู้พัฒนาภายนอกสามารถนำไปต่อยอดในงานวิจัยหรือเครื่องมืออื่น ๆ
        ได้ โดยไม่ต้องสร้าง pipeline เก็บข้อมูลใหม่.
\end{enumerate}

\section{โครงสร้างของวิทยานิพนธ์}
\label{sec:structure}

วิทยานิพนธ์ฉบับนี้ประกอบด้วย 5 บทหลักและภาคผนวก 3 ส่วน ดังนี้

\begin{itemize}[leftmargin=2em]
  \item \textbf{บทที่ 1 บทนำ} --- กล่าวถึงที่มาและความสำคัญของงานวิจัย
        ปัญหา วัตถุประสงค์ ขอบเขต และโครงสร้างของรายงาน.
  \item \textbf{บทที่ 2 การทบทวนวรรณกรรม} --- ทบทวนทฤษฎีพื้นฐานเกี่ยวกับ
        การคำนวณต้นทุนงานก่อสร้าง หลักการคำนวณต้นทุนต่อหน่วย เทคโนโลยี
        การดึงข้อมูล และงานวิจัยที่เกี่ยวข้อง.
  \item \textbf{บทที่ 3 วิธีการดำเนินงาน} --- อธิบายสถาปัตยกรรมระบบ
        การออกแบบฐานข้อมูล กลยุทธ์การดึงข้อมูล กระบวนการ canonicalization
        การคำนวณ BOQ และการวิเคราะห์เชิงสถิติ.
  \item \textbf{บทที่ 4 ผลการศึกษา} --- นำเสนอผลการทำงานของระบบ ผลการ
        ดึงข้อมูล ผลการเปรียบเทียบราคาข้ามแหล่ง ผลการวิเคราะห์แนวโน้ม
        และตัวอย่างผลการคำนวณ BOQ.
  \item \textbf{บทที่ 5 สรุปและข้อเสนอแนะ} --- สรุปผลการวิจัย ความ
        มีส่วนร่วมเชิงวิชาการ ข้อจำกัด และข้อเสนอแนะสำหรับงานวิจัย
        ในอนาคต.
  \item \textbf{ภาคผนวก ก} ภาพหน้าจอระบบ
  \item \textbf{ภาคผนวก ข} ตัวอย่างโค้ดสำคัญ
  \item \textbf{ภาคผนวก ค} ข้อมูลดิบที่ใช้ในระบบ
\end{itemize}
